\documentclass[11pt]{article}
\usepackage[margin=0.85in]{geometry}
\usepackage{amsmath,amssymb,amsfonts}
\usepackage{graphicx}
\usepackage{booktabs}
\usepackage{hyperref}
\usepackage{xcolor}
\usepackage{enumitem}
\usepackage{microtype}

\graphicspath{{./}}

\setlength{\parskip}{0.5em}
\setlength{\parindent}{0pt}

\title{\large\textbf{Optimizer Trajectories Mislead:\\
       Gradient Decomposition Recovers Feature Directions in Multitask Training}}
\author{Yongzhong Xu \quad\small\texttt{abbyxu@gmail.com}}
\date{\small\today}

\begin{document}
\maketitle
\vspace{-1.5em}

\noindent\rule{\linewidth}{0.4pt}

\paragraph{Core idea.}
Most interpretability work that analyzes training dynamics assumes:
\emph{directions extracted from optimizer trajectories reflect
feature formation}. This is wrong in a precise and fixable way.

\paragraph{Setup.}
\begin{itemize}[leftmargin=1.2em,itemsep=0.2em,topsep=0.2em]
  \item \textbf{Task:} grokking on modular arithmetic
        ($a{+}b$, $a{-}b$, $a{\cdot}b$, $a^2{+}b^2 \bmod 97$).
  \item \textbf{Probe:} the Linear Centroids Hypothesis (LCH)
        of Walker et al.\ \cite{lch_short} --- features as linear
        directions of input-space centroids
        $\mu_x = \nabla_x \ell_y(x)$, with sensitivity ratio
        $R_k(t)$ = centroid response along direction $v_k$ vs.\
        a random direction.
  \item \textbf{Compare two objects:}
    \begin{itemize}[leftmargin=1.2em,itemsep=0pt]
      \item Update SED: SVD over rolling AdamW updates $\Delta\theta_t$.
      \item Gradient SED: SVD over rolling gradients
            $g(t) = \nabla L|_{\theta_t}$.
    \end{itemize}
\end{itemize}

\paragraph{Result 1 --- Single-task.}
\begin{itemize}[leftmargin=1.2em,itemsep=0.2em,topsep=0.2em]
  \item Update SED: weak coupling, $\bar R_k \in [3, 9]$.
  \item Gradient SED: strong coupling, $\bar R_k \in [100, 330]$
        across 4 ops and 3 seeds.
\end{itemize}
But a $W$ ablation (window size) reveals: $R_1{=}648$ at $W{=}1$,
decaying monotonically to $140$ at $W{=}40$. The signal is
dominated by instantaneous gradient alignment, which is
near-tautological --- $\nabla_\theta L$ and $\mu_x$ are built from
the same Jacobians of the network. \emph{Rolling-window SVD acts as
denoising, not as the source of the coupling.}

\paragraph{Result 2 --- Multitask (the load-bearing finding).}
A multitask transformer with shared encoder and 4 op-specific heads,
trained on all four ops jointly:
\begin{itemize}[leftmargin=1.2em,itemsep=0.2em,topsep=0.2em]
  \item Update SED \emph{collapses}: $\bar R_k \le 1$ across all 4
        ops --- worse than random.
  \item Per-task gradient SED \emph{recovers signal}:
        $\bar R_k \in [20, 45]$, consistent across all 4 ops.
\end{itemize}

The aggregated AdamW update mixes (i) momentum, (ii) adaptive
scaling, (iii) weight decay, and (iv) inter-task gradient
cancellation, destroying alignment with per-task feature
directions. The fix: decompose gradients by task, then SVD per
task. This requires only one extra backward pass per task per
measurement checkpoint --- cheap relative to training.

\paragraph{Result 3 --- Causal intervention.}
Constraining the AdamW attention update to a rank-3 subspace,
3 seeds, addition and multiplication (15+15 = 30 runs):
\begin{center}\small
\begin{tabular}{lcc}
\toprule
mode (gradient-projected) & speedup over control \\
\midrule
remove SED rank-3        & $1.02\times$ \quad (no effect) \\
keep only SED rank-3     & $\mathbf{2.20\times}$ \\
remove random rank-3     & $1.00\times$ \quad (no effect) \\
keep only random rank-3  & $\mathbf{2.36\times}$ \\
\bottomrule
\end{tabular}
\end{center}
The acceleration is \textbf{rank-specific, not direction-specific}.
Any rank-3 keep yields $\approx 2.3\times$ speedup; SED-vs-random
is within seed variation. So the SED basis is \emph{diagnostic} of
where centroid feature formation is concentrated, but \emph{not
causally privileged}: the rank constraint does the work.

\paragraph{Takeaways.}
\begin{enumerate}[leftmargin=1.2em,itemsep=0.2em,topsep=0.2em]
  \item Optimizer trajectories $\ne$ feature directions.
  \item Gradient structure is the right object.
  \item Multitask requires per-task gradient decomposition.
  \item Low-rank structure matters more than specific directions
        for whether grokking happens; specific directions matter
        for measuring \emph{where} it happens.
\end{enumerate}

\paragraph{One-line punchline.}
\emph{If you want to understand what a network is learning,
look at gradients --- not what AdamW did to them.}

\paragraph{Why this matters.}
\begin{itemize}[leftmargin=1.2em,itemsep=0pt,topsep=0pt]
  \item Explains apparent failures of trajectory PCA / SVD
        diagnostics in multitask settings.
  \item Connects to multitask gradient-interference literature.
  \item Suggests a candidate diagnostic for monitoring per-task
        feature formation in instruction-tuning, MoE, and other
        shared-encoder pipelines, where gradient interference is
        pervasive --- though we have only verified the claim on a
        toy modular-arithmetic setting and large-scale validation
        is open.
\end{itemize}

\noindent\rule{\linewidth}{0.4pt}
{\footnotesize Full paper, code, and per-checkpoint data caches:
\url{https://github.com/skydancerosel/grokking-integrability}}

\begin{thebibliography}{1}
\bibitem{lch_short}
Walker, Humayun, Balestriero, Baraniuk. \emph{The Linear Centroids
Hypothesis: How Deep Network Features Represent Data.}
arXiv:2604.11962 (2026).
\end{thebibliography}

\end{document}
